(1) 一维稳态无内热源时热流密度 $j$ 为常数，各区温差满足 $\Delta T_1=jd_1/k_1$、$\Delta T_c=j\mathcal R_c$、$\Delta T_2=jd_2/k_2$。串联相加得 $j=(T_h-T_c)/(d_1/k_1+\mathcal R_c+d_2/k_2)$。从热端到界面左侧 $T_{1I}=T_h-jd_1/k_1$，界面右侧 $T_{2I}=T_{1I}-j\mathcal R_c$，两者确实存在正的温度跃变。
(2) 在板内局部熵产生密度为 $k(\nabla T)^2/T^2$，接触面熵产生率为 $j(1/T_{2I}-1/T_{1I})$。积分各板并利用 $j=-k\,dT/dx$，总单位面积熵产生率望远镜式相加为 $\dot s_{\rm gen}=j(1/T_c-1/T_h)>0$。虽然接触热阻使局部温度不连续，第二定律仍要求各部分对应热流均从高温流向低温。
(3) 若把界面附近小系统的面热容取为 $\mathcal C$，采用两端热导 $G_h,G_c$ 的集总模型，能量平衡为 $\mathcal C\,dT_I/dt=G_h(T_h-T_I)-G_c(T_I-T_c)$。稳态温度 $T_{I*}=(G_hT_h+G_cT_c)/(G_h+G_c)$；通解 $T_I(t)=T_{I*}+[T_I(0)-T_{I*}]e^{-t/\tau}$，$\tau=\mathcal C/(G_h+G_c)$，可由直接代回微分方程复核。
(4) 当 $k_1\to\infty$ 时第一层温降消失，热流仍受 $\mathcal R_c+d_2/k_2$ 限制；当 $\mathcal R_c\to\infty$ 时 $j\to0$，两个板分别接近各自边界温度而界面几乎绝热。提高已经很小的 $d_1/k_1$ 对串联总热阻改变极少，设计上应优先降低主导接触热阻。
